\chapter{\IfLanguageName{dutch}{Stand van zaken}{State of the art}}%
\label{ch:stand-van-zaken}

Dit hoofdstuk schetst het theoretisch kader voor het evalueren van kredietrisico's en de rol van datatechnologie daarin. Om de centrale onderzoeksvraag te kunnen beantwoorden, analyseert de literatuurstudie eerst de traditionele methodieken rond kredietrisicoanalyse en de beperkingen van de huidige softwareoplossingen. Vervolgens verschuift de focus naar de technologische context, waarbij de mogelijkheden van cloud-datawarehouses en machine learning-algoritmen (ML) binnen de financiële sector (FinTech) worden belicht.

\section{Traditionele Risicoanalyse en Bestaande Software}

De financiële gezondheid van een onderneming evalueert men traditioneel aan de hand van ratio-analyse. Analisten berekenen indicatoren zoals liquiditeit, solvabiliteit en rentabiliteit op basis van de neergelegde jaarrekeningen \autocite{Ooghe2021}. Modellen, zoals de Altman Z-score, proberen op basis van deze ratio's de kans op een eventueel faillissement te voorspellen.

Om dit proces efficiënter te maken voor de analist, biedt de markt diverse softwaretoepassingen aan, waaronder Companyweb, Belfirst en OpenTheBox. Deze platformen slagen erin bedrijfsstructuren te visualiseren en geven toegang tot sectorgemiddelde ratio's. Hoewel deze tools uitstekend presteren voor het raadplegen van individuele klantdossiers, missen ze vaak de flexibiliteit om grootschalige, op maat gemaakte predictieve modellering uit te voeren op ruwe data. Gebruikers zijn gebonden aan de visualisaties en berekeningen die de softwareleverancier voorziet, wat exploratieve analyse op grotere schaal belemmert.

\section{Big Data en Cloud Computing in FinTech}

De opkomst van 'FinTech' brengt een verschuiving teweeg van statische analyses naar dynamische, datagedreven toepassingen. Binnen dit domein zijn cloud-native datawarehouses cruciaal. Architecturen zoals Google BigQuery vertonen superieure prestaties bij het verwerken van datasets met miljoenen rijen in vergelijking met traditionele, on-premise servers \autocite{Gupta2024}. BigQuery typeert zich door een serverless opbouw en laat toe om complexe data analyses uit te voeren op een kostenefficiënte en uiterst performante manier \autocite{ActiveViam2025}.

% (Deel 1 en 2 blijven behouden zoals je had, hier is de correctie vanaf sectie 3) ...

\section{Machine Learning voor Kredietrisico}

Naast de opslag en verwerking van Big Data, besteedt de literatuur veel aandacht aan de inzet van Machine Learning voor het voorspellen van kredietrisico's. Recente studies tonen aan dat moderne datagedreven technieken aanzienlijk beter kunnen presteren dan traditionele rekenmethodes. Volgens \textcite{Moscato2021} bieden vooral geavanceerde algoritmen, zoals ensemble learning, een hogere accuraatheid bij \textit{credit scoring}. De integratie van zulke ML-technieken, gecombineerd met de rekenkracht van cloud-datawarehouses en gestructureerde datasets zoals die van de Nationale Bank van België, kan hierdoor nieuwe perspectieven bieden voor een proactiever en nauwkeuriger acceptatiebeleid in de kredietverlening.

Een kritieke factor bij de ontwikkeling van financieel voorspellende machine learning-modellen is de kwaliteit en structuur van de trainingsdataset. Het gebruik van één generiek model om de financiële ratio's van alle ondernemingen te voorspellen, levert doorgaans onnauwkeurige resultaten op wegens de grote onderlinge absolute verschillen en marktdynamieken. Er moet bijgevolg per financieel zinvolle cluster een specifiek model worden ontwikkeld om de voorspellingen binnen de cluster (intra-cluster) te optimaliseren.

In de academische literatuur worden doorgaans twee methoden onderscheiden om deze clusters op te bouwen: een handmatige segmentatie op basis van vooraf gedefinieerde parameters, of een algoritmische clustering (zoals K-means). Binnen de traditionele segmentatie zijn de voornaamste determinanten voor het groeperen van ondernemingen reeds afgebakend. Ten eerste vormt de grootte van de onderneming een fundamentele variabele; \textcite{cultrera2016} stellen immers dat grotere bedrijven (die het volledig schema hanteren) profiteren van schaalvoordelen, een vlottere toegang hebben tot bankkredieten en een totaal ander werkkapitaalbeheer hanteren dan micro-ondernemingen. Een micro-entiteit wordt bovendien vaak gestuurd door de persoonlijke financiën van de zaakvoerder (bijv. via de rekening-courant), wat de solvabiliteitsratio's scheeftrekt vergeleken met een grote NV. Een algoritme dat deze profielen door elkaar haalt, leert patronen aan die in realiteit niet bestaan.

Ten tweede is de leeftijd (de levenscyclus) van de onderneming een bepalende factor, waarbij \textcite{cho2020} aantonen dat het voorspellend vermogen toeneemt wanneer modellen worden afgestemd op de maturiteit van de entiteit. Tot slot is een sectoriële afbakening noodzakelijk, aangezien financiële structuren sterk fluctueren naargelang de kernactiviteit en het gebruik van deze afbakening een significante en bewezen impact heeft op de accuraatheid van insolventiemodellen \autocite{tian2017}.

Aan de andere kant toont recent onderzoek aan dat algoritmische clusteringtechnieken in staat zijn om meer homogene groepen te vormen dan deze traditionele opsplitsingsmethoden. Door hun datagedreven, \textit{unsupervised} karakter kunnen dergelijke algoritmen complexe, onderliggende patronen in de data identificeren, zonder daarbij afhankelijk te zijn van rigide, vooraf bepaalde labels \autocite{hou2016financial, gulal2023predicting, zhao2024boosting}.toont aan dat clusteringtechnieken in staat zijn om meer homogene groepen te vormen dan traditionele opsplitsingsmethoden. Door hun datagedreven, \textit{unsupervised} karakter kunnen deze algoritmen complexe, onderliggende patronen in de data identificeren, zonder daarbij afhankelijk te zijn van rigide, vooraf bepaalde labels \cite{hou2016financial, gulal2023predicting, zhao2024boosting}.